% Examen ALEA du 9 janvier 2026, exercice 3.
% Le PDF imprime x\sqrt{2\pi\sigma} au dénominateur : coquille corrigée
% en x\sigma\sqrt{2\pi}, nécessaire pour obtenir une densité de probabilité.
\uuid{lu9y}
\titre{Loi log-normale et produits de variables aléatoires}
\niveau{L2}
\module{Probabilité}
\chapitre{Probabilité continue}
\sousChapitre{Loi normale}
\theme{Loi log-normale, changement de variable, produit}
\auteur{}
\datecreate{2026-10-01}
\organisation{AMSCC}
\difficulte{}
\contenu{
\texte{La loi log-normale de paramètres $\mu\in\mathbb{R}$ et $\sigma>0$, notée $\mathcal{LN}(\mu,\sigma)$, est la loi de densité
\[f(x)=\frac{1}{x\sigma\sqrt{2\pi}}\exp\!\left(-\frac{(\ln x-\mu)^2}{2\sigma^2}\right)\mathbf{1}_{]0,+\infty[}(x).\]}
\begin{enumerate}
  \item \question{Soit $X$ une variable aléatoire à valeurs strictement positives. Montrer que $Y=\ln X$ suit la loi normale $\mathcal{N}(\mu,\sigma)$ si et seulement si $X$ suit la loi $\mathcal{LN}(\mu,\sigma)$.}
  \item \question{Montrer que, si $X$ suit la loi $\mathcal{LN}(\mu,\sigma)$, alors $1/X$ suit une loi log-normale dont on précisera les paramètres.}
  \item \question{Soit $(X_i)_{1\leq i\leq n}$ une famille de variables aléatoires mutuellement indépendantes avec $X_i\sim\mathcal{LN}(\mu_i,\sigma_i)$. Déterminer la loi de $Z=\prod_{i=1}^{n}X_i$.}
  \item \question{On suppose que $X\sim\mathcal{LN}(5,1{,}5)$. Calculer $\mathbb{P}(50\leq X\leq 400)$.}
  \item \question{Soit $X\sim\mathcal{LN}(0,1)$. Montrer que
  \[\mathbb{E}[X]=\frac{1}{\sqrt{2\pi}}\int_{\mathbb{R}} e^{y-y^2/2}\,\mathrm{d}y=\exp\!\left(\frac12\right).\]}
\end{enumerate}
}
